\documentclass[12pt]{article}

% Pachete necesare pentru codificarea caracterelor și diacritice
\usepackage[utf8]{inputenc}
\usepackage[romanian]{babel}

% Pachete matematice avansate
\usepackage{amsmath}
\usepackage{amsfonts}
\usepackage{amssymb}
\usepackage{amsthm}
\usepackage{booktabs} % Pentru tabele elegante

% Setarea marginilor paginii
\usepackage[margin=1in]{geometry}

\title{Seminar 2: Criptografie -- Rezolvare Teme}
\author{}
\date{} % Elimină complet data și spațiul ei de sub titlu

\begin{document}

\maketitle


\subsection*{Cerință}
\textit{Determinați cel mai mare divizor comun al lui $32335$ și $53331$ folosind algoritmul lui Euclid extins și găsiți coeficienții Bezout.}


\subsection*{Algoritmul lui Euclid Extins (Determinarea coeficienților)}
Pentru a găsi coeficienții întregi $u$ și $v$ care satisfac identitatea lui Bézout:
\[
29 = u \cdot 53331 + v \cdot 32335
\]
urmărim evoluția vectorilor de coeficienți $x_i = (u_i, v_i)$ asociați fiecărui rest curent, folosind relația de recurență:
\[
x_{i} = x_{i-2} - q_{i-2} \cdot x_{i-1}
\]
unde $q_{i-2}$ este câtul obținut la pasul respectiv.

\subsubsection*{Inițializare:}
\begin{itemize}
    \item Pentru $a = 53331$: \quad $x_a = (1, 0)$
    \item Pentru $b = 32335$: \quad $x_b = (0, 1)$
\end{itemize}

\subsubsection*{Calculul pas cu pas:}
\begin{itemize}
    \item \textbf{Pasul 1} (Câtul $q_0 = 1$): \\
    $x_{20996} = (1, 0) - 1 \times (0, 1) = (1, -1)$
    
    \item \textbf{Pasul 2} (Câtul $q_1 = 1$): \\
    $x_{11339} = (0, 1) - 1 \times (1, -1) = (-1, 2)$
    
    \item \textbf{Pasul 3} (Câtul $q_2 = 1$): \\
    $x_{9657} = (1, -1) - 1 \times (-1, 2) = (2, -3)$
    
    \item \textbf{Pasul 4} (Câtul $q_3 = 1$): \\
    $x_{1682} = (-1, 2) - 1 \times (2, -3) = (-3, 5)$
    
    \item \textbf{Pasul 5} (Câtul $q_4 = 5$): \\
    $x_{1247} = (2, -3) - 5 \times (-3, 5) = (2 - (-15), -3 - 25) = (17, -28)$
    
    \item \textbf{Pasul 6} (Câtul $q_5 = 1$): \\
    $x_{435} = (-3, 5) - 1 \times (17, -28) = (-20, 33)$
    
    \item \textbf{Pasul 7} (Câtul $q_6 = 2$): \\
    $x_{377} = (17, -28) - 2 \times (-20, 33) = (17 - (-40), -28 - 66) = (57, -94)$
    
    \item \textbf{Pasul 8} (Câtul $q_7 = 1$): \\
    $x_{58} = (-20, 33) - 1 \times (57, -94) = (-77, 127)$
    
    \item \textbf{Pasul 9} (Câtul $q_8 = 6$): \\
    $x_{29} = (57, -94) - 6 \times (-77, 127) = (57 - (-462), -94 - 762) = (519, -856)$
\end{itemize}

\subsection*{ Centralizarea Datelor}

Toate iterațiile pot fi sintetizate în următorul tabel de monitorizare a stării algoritmului:

\begin{center}
\begin{tabular}{crrrr}
\toprule
\textbf{Pasul ($i$)} & \textbf{Restul ($r_i$)} & \textbf{Câtul ($q_i$)} & \textbf{Coeficientul $u$} & \textbf{Coeficientul $v$} \\
\midrule
- & 53331 & - & 1 & 0 \\
- & 32335 & 1 & 0 & 1 \\
1 & 20996 & 1 & 1 & -1 \\
2 & 11339 & 1 & -1 & 2 \\
3 & 9657 & 1 & 2 & -3 \\
4 & 1682 & 5 & -3 & 5 \\
5 & 1247 & 1 & 17 & -28 \\
6 & 435 & 2 & -20 & 33 \\
7 & 377 & 1 & 57 & -94 \\
8 & 58 & 6 & -77 & 127 \\
\textbf{9} & \textbf{29} & 2 & \textbf{519} & \textbf{-856} \\
10 & 0 & - & - & - \\
\bottomrule
\end{tabular}
\end{center}

\subsection*{Concluzie }
Coeficienții Bezout determinați sunt $\mathbf{u = 519}$ și $\mathbf{v = -856}$. Identitatea lui Bézout se scrie:
\[
29 = 519 \cdot 53331 + (-856) \cdot 32335
\]


\end{document}